\section{Related Work}
\subsection{Datasets and Modular OMR}
OMR research has historically decomposed a page into staff detection, symbol recognition, and notation reconstruction \cite{calvozaragoza2021understanding,rebelo2012optical}. Dataset design reflects this decomposition. MUSCIMA++ provides handwritten symbols and explicit pairwise relations, enabling notation-graph experiments \cite{hajic2017muscima}. DeepScores and DeepScoresV2 provide dense synthetic annotations for tiny printed objects, instance segmentation, and detection \cite{tuggener2018deepscores,tuggener2020deepscoresv2}. Camera-PrIMuS supports end-to-end recognition under realistic image degradation but is primarily monophonic \cite{calvozaragoza2018primus}. RealScores and recent instance-segmentation studies further emphasize that real-world score images differ materially from synthetic renders \cite{shatri2023realworld,shatri2024segmentation}.

Object-centric OMR benefits from detectors that preserve small, dense instances. DETR introduced set prediction with bipartite matching \cite{carion2020detr}; RT-DETRv2, D-FINE, and DEIM improve multiscale sampling, localization, and convergence \cite{lv2024rtdetrv2,peng2025dfine,huang2025deim}. Promptable segmentation offers a complementary route from boxes to boundaries. SAM established zero-shot promptable segmentation \cite{kirillov2023segment}, while SAM2 extended the design and improved image efficiency \cite{ravi2025sam2}. RIGOR-OMR uses these models as reusable components, not as evidence that natural-image priors alone understand notation semantics. DINOv2 features support retrieval and crop-level auxiliary classification \cite{oquab2023dinov2}, but all final claims are based on end-output metrics rather than embedding similarity.

\subsection{End-to-End Pianoform Recognition}
GrandStaff introduced a large pianoform corpus and a sequence-recognition formulation over kern-like encodings \cite{riosvila2023grandstaff}. Sheet Music Transformer extended image-to-sequence OMR beyond monophonic transcription \cite{riosvila2024smt}, and full-page variants combined convolutional encoders, autoregressive decoders, synthetic curriculum learning, and target-domain fine-tuning \cite{riosvila2026fullpage}. TrOMR likewise targets polyphonic recognition with Transformer-based decoding \cite{li2023tromr}. OLiMPiC introduced LMX, aligned MusicXML-compatible output, scanned evaluation systems, and a tree-aware metric \cite{mayer2024practical}. These systems demonstrate strong in-domain accuracy, but their results also motivate a controlled view of cross-domain behavior.

\subsection{Transfer and Domain Shift}
Transfer learning studies distinguish reusable lower-level representations from source-specific higher layers \cite{pan2010transfer,yosinski2014transferable}. Domain-adversarial training seeks features that cannot discriminate source from target domains \cite{ganin2016domain}, and domain-generalization work organizes interventions into data, representation, and learning strategies \cite{wang2021generalizing}. RIGOR-OMR studies a narrower and practically common setting: an OMR system must operate on an unseen collection without target annotations, while appearance, layout, and output representation may shift simultaneously.
